% ---------------------------------------------------------------------
% GENERATED by tools/make_api_appendix.py: do not edit by hand,
% run the script again after changing the code.
% ---------------------------------------------------------------------
\chapter{Reference of classes, functions and methods}
\label{app:api}

This appendix lists, module by module, every class, function and method of the package, with the signature and a one-sentence description taken from the docstrings. Methods are written as \code{Class.method(arguments)}. Names beginning with an underscore are internal helpers, listed for completeness. The full explanation of each item is in its docstring and in the chapters of this manual.

\section{\texorpdfstring{\texttt{beam.py}}{beam.py}: Incident beam}
\label{app:api:beam}
Where do the particles of a rectangular incident beam start?

\begin{longtable}{@{}p{1.75cm}p{6.35cm}p{6.45cm}@{}}
\toprule \textbf{Kind} & \textbf{Name and arguments} & \textbf{What it does} \\ \midrule \endhead
\small function & \raggedright\arraybackslash\small \texttt{\_\allowbreak{}sample\_\allowbreak{}1d(length, profile, n, rng, n\_\allowbreak{}grid=2001)} & \raggedright\arraybackslash\small Draw \texttt{n} random positions in [-length/2, +length/2] distributed according to \texttt{profile}. \emph{(internal)} \\
\small function & \raggedright\arraybackslash\small \texttt{sample\_\allowbreak{}beam\_\allowbreak{}point(width, height, H=None, V=None, n=1, rng=None)} & \raggedright\arraybackslash\small Pick the (x, y) point(s) in the beam rectangle where particle(s) start. \\
\small function & \raggedright\arraybackslash\small \texttt{start\_\allowbreak{}position(width, height, H=None, V=None, z0=0.\allowbreak{}0, rng=None)} & \raggedright\arraybackslash\small Convenience wrapper: one particle, returned as a 3D point (x, y, z0) together with its direction of flight (0, 0, 1), i.e. along +z. \\
\bottomrule
\end{longtable}

\section{\texorpdfstring{\texttt{sample\_base.py}}{sample_base.py}: Common ingredients of all samples}
\label{app:api:sample_base}
Ingredients shared by ALL sample geometries of MCAMS.

\begin{longtable}{@{}p{1.75cm}p{6.35cm}p{6.45cm}@{}}
\toprule \textbf{Kind} & \textbf{Name and arguments} & \textbf{What it does} \\ \midrule \endhead
\small dataclass & \raggedright\arraybackslash\small \texttt{Layer(name: str, thickness: float, mu: float, scatter\_\allowbreak{}fraction: float = 1.\allowbreak{}0, fill\_\allowbreak{}height: Optional[float] = None, fullness: Optional[float] = None, bottom: Optional[float] = None, height: Optional[float] = None)} & \raggedright\arraybackslash\small One layer of the sample. \\
\small property & \raggedright\arraybackslash\small \texttt{Layer.\allowbreak{}is\_\allowbreak{}partial} & \raggedright\arraybackslash\small True if the layer was declared with a limited vertical extent (bottom, height, fill\_height, or fullness < 1). \\
\small function & \raggedright\arraybackslash\small \texttt{make\_\allowbreak{}layer(item)} & \raggedright\arraybackslash\small Turn one item of a list of layers into a \texttt{Layer}. \\
\small function & \raggedright\arraybackslash\small \texttt{scattered\_\allowbreak{}direction(two\_\allowbreak{}theta\_\allowbreak{}deg, phi\_\allowbreak{}deg=0.\allowbreak{}0)} & \raggedright\arraybackslash\small Unit vector (LAB frame) of a particle scattered by \texttt{two\_\allowbreak{}theta\_\allowbreak{}deg} from the beam direction +z. \\
\small class & \raggedright\arraybackslash\small \texttt{LayeredSample(ABC)} & \raggedright\arraybackslash\small Abstract base class of the layered samples (slab, cylinder, ...). \\
\small init & \raggedright\arraybackslash\small \texttt{LayeredSample.\allowbreak{}\_\allowbreak{}\_\allowbreak{}init\_\allowbreak{}\_\allowbreak{}(layers, height=None, y\_\allowbreak{}bottom=None, beam\_\allowbreak{}height=None)} & \raggedright\arraybackslash\small Store the layer data and the vertical extent of the sample. \\
\small method & \raggedright\arraybackslash\small \texttt{LayeredSample.\allowbreak{}\_\allowbreak{}material\_\allowbreak{}range\_\allowbreak{}of(layer)} & \raggedright\arraybackslash\small (bottom, top) heights y (mm) of the material of one layer, limited to the sample [y\_bottom, y\_top]. \emph{(internal)} \\
\small method & \raggedright\arraybackslash\small \texttt{LayeredSample.\allowbreak{}beam\_\allowbreak{}fullness(beam\_\allowbreak{}height=None)} & \raggedright\arraybackslash\small Fraction of the beam height covered by the material of each layer. \\
\small property & \raggedright\arraybackslash\small \texttt{LayeredSample.\allowbreak{}has\_\allowbreak{}partial\_\allowbreak{}layers} & \raggedright\arraybackslash\small True if at least one layer does not fill the whole sample height. \\
\small method & \raggedright\arraybackslash\small \texttt{LayeredSample.\allowbreak{}is\_\allowbreak{}filled(layer\_\allowbreak{}index, y, tol=GEOMETRY\_\allowbreak{}TOL)} & \raggedright\arraybackslash\small True where the height y (mm) is inside the MATERIAL of the given layer(s), i.e. between its bottom and its top. \\
\small method & \raggedright\arraybackslash\small \texttt{LayeredSample.\allowbreak{}\_\allowbreak{}fill\_\allowbreak{}text(i)} & \raggedright\arraybackslash\small Short text with the vertical extent of layer i, used by \texttt{summary}: "full", or "bottom to top" (mm) and, when the beam height is known, the fullness in parentheses. \emph{(internal)} \\
\small property & \raggedright\arraybackslash\small \texttt{LayeredSample.\allowbreak{}n\_\allowbreak{}layers} & \raggedright\arraybackslash\small Number of layers. \\
\small method & \raggedright\arraybackslash\small \texttt{LayeredSample.\allowbreak{}transmission(path\_\allowbreak{}lengths)} & \raggedright\arraybackslash\small Fraction of particles surviving the given paths: exp(-sum mu\_i*L\_i). \\
\small abstract & \raggedright\arraybackslash\small \texttt{LayeredSample.\allowbreak{}summary()} & \raggedright\arraybackslash\small Return a text table describing the sample. \\
\small abstract & \raggedright\arraybackslash\small \texttt{LayeredSample.\allowbreak{}to\_\allowbreak{}sample\_\allowbreak{}frame(vectors)} & \raggedright\arraybackslash\small Convert vectors (directions) from the lab frame to the sample frame. \\
\small abstract & \raggedright\arraybackslash\small \texttt{LayeredSample.\allowbreak{}to\_\allowbreak{}lab\_\allowbreak{}frame(vectors)} & \raggedright\arraybackslash\small Inverse of \texttt{to\_\allowbreak{}sample\_\allowbreak{}frame}. \\
\small abstract & \raggedright\arraybackslash\small \texttt{LayeredSample.\allowbreak{}contains(points, tol=1e-09)} & \raggedright\arraybackslash\small True for each point (SAMPLE frame, shape (N, 3)) inside the volume of the sample (including the empty part of partially filled layers). \\
\small abstract & \raggedright\arraybackslash\small \texttt{LayeredSample.\allowbreak{}layer\_\allowbreak{}of\_\allowbreak{}points(points)} & \raggedright\arraybackslash\small Index of the layer whose MATERIAL contains each point (SAMPLE frame, shape (N, 3)); -1 for points outside the sample or in the empty part above a partially filled layer. \\
\small abstract & \raggedright\arraybackslash\small \texttt{LayeredSample.\allowbreak{}outgoing\_\allowbreak{}path\_\allowbreak{}lengths(points, directions)} & \raggedright\arraybackslash\small Distance travelled in the material of each layer by particles that start at \texttt{points} and fly along \texttt{directions} (both in the SAMPLE frame) until they leave the sample. \\
\small abstract & \raggedright\arraybackslash\small \texttt{LayeredSample.\allowbreak{}next\_\allowbreak{}interaction(points, directions, rng=None)} & \raggedright\arraybackslash\small For particles flying from \texttt{points} along \texttt{directions} (SAMPLE frame, inside the sample): where is their next interaction? Returns (new\_points (N, 3), interacts (N,) boolean). \\
\small abstract & \raggedright\arraybackslash\small \texttt{LayeredSample.\allowbreak{}sample\_\allowbreak{}interaction\_\allowbreak{}points(x, y, rng=None)} & \raggedright\arraybackslash\small First interaction of beam particles. \\
\bottomrule
\end{longtable}

\section{\texorpdfstring{\texttt{slab.py}}{slab.py}: Slab geometry}
\label{app:api:slab}
SLAB geometry of MCAMS: a rectangular sample made of flat layers.

\begin{longtable}{@{}p{1.75cm}p{6.35cm}p{6.45cm}@{}}
\toprule \textbf{Kind} & \textbf{Name and arguments} & \textbf{What it does} \\ \midrule \endhead
\small function & \raggedright\arraybackslash\small \texttt{\_\allowbreak{}slice\_\allowbreak{}interval(p, d, lo, hi)} & \raggedright\arraybackslash\small Distances t at which rays p + t*d are between two parallel planes. \emph{(internal)} \\
\small class & \raggedright\arraybackslash\small \texttt{SlabSample(LayeredSample)} & \raggedright\arraybackslash\small Rectangular multilayer slab sample. \\
\small init & \raggedright\arraybackslash\small \texttt{SlabSample.\allowbreak{}\_\allowbreak{}\_\allowbreak{}init\_\allowbreak{}\_\allowbreak{}(width, height, layers, z\_\allowbreak{}front=0.\allowbreak{}0, plane\_\allowbreak{}angle\_\allowbreak{}deg=90.\allowbreak{}0, y\_\allowbreak{}bottom=None, beam\_\allowbreak{}height=None)} & \raggedright\arraybackslash\small Constructor: creates the object (see the description of the class). \\
\small property & \raggedright\arraybackslash\small \texttt{SlabSample.\allowbreak{}total\_\allowbreak{}thickness} & \raggedright\arraybackslash\small Total thickness of the stack along the normal (mm). \\
\small method & \raggedright\arraybackslash\small \texttt{SlabSample.\allowbreak{}summary()} & \raggedright\arraybackslash\small Return a small text table describing the sample. \\
\small method & \raggedright\arraybackslash\small \texttt{SlabSample.\allowbreak{}with\_\allowbreak{}plane\_\allowbreak{}angle(plane\_\allowbreak{}angle\_\allowbreak{}deg)} & \raggedright\arraybackslash\small Return a copy of this sample with another orientation. \\
\small method & \raggedright\arraybackslash\small \texttt{SlabSample.\allowbreak{}to\_\allowbreak{}sample\_\allowbreak{}frame(vectors)} & \raggedright\arraybackslash\small Convert vectors (e.g. directions) from the lab frame to the sample frame. \\
\small method & \raggedright\arraybackslash\small \texttt{SlabSample.\allowbreak{}to\_\allowbreak{}lab\_\allowbreak{}frame(vectors)} & \raggedright\arraybackslash\small Inverse of \texttt{to\_\allowbreak{}sample\_\allowbreak{}frame}: sample-frame vectors -> lab frame. \\
\small method & \raggedright\arraybackslash\small \texttt{SlabSample.\allowbreak{}contains(points, tol=1e-09)} & \raggedright\arraybackslash\small True for every point (x', y', z') that lies inside the box of the sample (the empty parts of partially filled layers included). \\
\small method & \raggedright\arraybackslash\small \texttt{SlabSample.\allowbreak{}layer\_\allowbreak{}index(z)} & \raggedright\arraybackslash\small Index of the layer containing depth z' (-1 if outside the sample). \\
\small method & \raggedright\arraybackslash\small \texttt{SlabSample.\allowbreak{}layer\_\allowbreak{}of\_\allowbreak{}points(points)} & \raggedright\arraybackslash\small Layer index of each point (x', y', z') of an (N, 3) array. \\
\small method & \raggedright\arraybackslash\small \texttt{SlabSample.\allowbreak{}\_\allowbreak{}section\_\allowbreak{}interval(p, d)} & \raggedright\arraybackslash\small Interval of t during which rays p + t*d are inside the OUTLINE of the sample (all layers together, empty parts included). \emph{(internal)} \\
\small method & \raggedright\arraybackslash\small \texttt{SlabSample.\allowbreak{}\_\allowbreak{}clip\_\allowbreak{}inside(points)} & \raggedright\arraybackslash\small Push points that rounding errors put just outside the sample back onto its surface (here: limit each coordinate to the box). \emph{(internal)} \\
\small method & \raggedright\arraybackslash\small \texttt{SlabSample.\allowbreak{}\_\allowbreak{}material\_\allowbreak{}intervals(points, directions)} & \raggedright\arraybackslash\small For rays p + t*d (t >= 0), the interval [t0\_i, t1\_i] spent in the MATERIAL of each layer i. \emph{(internal)} \\
\small static & \raggedright\arraybackslash\small \texttt{SlabSample.\allowbreak{}\_\allowbreak{}unit(points, directions)} & \raggedright\arraybackslash\small Points and unit directions as (N, 3) float arrays (one direction is repeated for all the points). \emph{(internal)} \\
\small method & \raggedright\arraybackslash\small \texttt{SlabSample.\allowbreak{}outgoing\_\allowbreak{}path\_\allowbreak{}lengths(points, directions)} & \raggedright\arraybackslash\small Path travelled in the material of each layer by a particle that starts at each point and flies along \texttt{directions} until it leaves the sample (through the front, back or any side face, including the top). \\
\small method & \raggedright\arraybackslash\small \texttt{SlabSample.\allowbreak{}incident\_\allowbreak{}path\_\allowbreak{}lengths(points)} & \raggedright\arraybackslash\small Path travelled in the material of each layer by the incident beam (direction +z of the lab frame) from where it enters the sample up to each point (sample frame), in mm, shape (N, n\_layers). \\
\small method & \raggedright\arraybackslash\small \texttt{SlabSample.\allowbreak{}\_\allowbreak{}first\_\allowbreak{}interaction(points, directions, rng)} & \raggedright\arraybackslash\small Find the first interaction of particles flying from \texttt{points} along \texttt{directions} (SAMPLE frame; the points may be inside the sample or upstream of it). \emph{(internal)} \\
\small method & \raggedright\arraybackslash\small \texttt{SlabSample.\allowbreak{}\_\allowbreak{}depth\_\allowbreak{}from\_\allowbreak{}tau(tau\_\allowbreak{}normal)} & \raggedright\arraybackslash\small Convert an optical depth measured along the layer normal into the depth z' where it is reached. \emph{(internal)} \\
\small method & \raggedright\arraybackslash\small \texttt{SlabSample.\allowbreak{}sample\_\allowbreak{}interaction\_\allowbreak{}depth(n=1, rng=None)} & \raggedright\arraybackslash\small Draw the depth z' (mm) at which \texttt{n} beam particles entering through the front face interact (Beer-Lambert law, oblique incidence taken into account). \\
\small method & \raggedright\arraybackslash\small \texttt{SlabSample.\allowbreak{}sample\_\allowbreak{}interaction\_\allowbreak{}points(x, y, rng=None)} & \raggedright\arraybackslash\small Find where beam particles interact inside the sample. \\
\small method & \raggedright\arraybackslash\small \texttt{SlabSample.\allowbreak{}next\_\allowbreak{}interaction(points, directions, rng=None)} & \raggedright\arraybackslash\small Follow particles that fly from \texttt{points} along \texttt{directions} (both in the SAMPLE frame, inside the sample) and find where each one has its next interaction, if it has one before leaving the sample. \\
\small method & \raggedright\arraybackslash\small \texttt{SlabSample.\allowbreak{}\_\allowbreak{}check\_\allowbreak{}points(points)} & \raggedright\arraybackslash\small Return points as an (N, 3) float array; raise if outside sample. \emph{(internal)} \\
\small function & \raggedright\arraybackslash\small \texttt{create\_\allowbreak{}slab\_\allowbreak{}sample(width, height, layers, z\_\allowbreak{}front=0.\allowbreak{}0, plane\_\allowbreak{}angle\_\allowbreak{}deg=90.\allowbreak{}0, y\_\allowbreak{}bottom=None, beam\_\allowbreak{}height=None)} & \raggedright\arraybackslash\small Build a layered rectangular (slab) sample. \\
\bottomrule
\end{longtable}

\section{\texorpdfstring{\texttt{disc.py}}{disc.py}: Disc geometry (version 2.3)}
\label{app:api:disc}
DISC geometry of MCAMS (version 2.3): a flat sample with a.

\begin{longtable}{@{}p{1.75cm}p{6.35cm}p{6.45cm}@{}}
\toprule \textbf{Kind} & \textbf{Name and arguments} & \textbf{What it does} \\ \midrule \endhead
\small function & \raggedright\arraybackslash\small \texttt{\_\allowbreak{}cylinder\_\allowbreak{}interval(px, py, dx, dy, radius)} & \raggedright\arraybackslash\small Distances t at which rays p + t*d are inside the infinite cylinder x\textasciicircum{}2 + y\textasciicircum{}2 <= radius\textasciicircum{}2 (axis along z). \emph{(internal)} \\
\small class & \raggedright\arraybackslash\small \texttt{DiscSample(SlabSample)} & \raggedright\arraybackslash\small Multilayer disc (circular plate) centred on the beam. \\
\small init & \raggedright\arraybackslash\small \texttt{DiscSample.\allowbreak{}\_\allowbreak{}\_\allowbreak{}init\_\allowbreak{}\_\allowbreak{}(diameter, layers, z\_\allowbreak{}front=0.\allowbreak{}0, plane\_\allowbreak{}angle\_\allowbreak{}deg=90.\allowbreak{}0, beam\_\allowbreak{}height=None)} & \raggedright\arraybackslash\small Constructor: creates the object (see the description of the class). \\
\small method & \raggedright\arraybackslash\small \texttt{DiscSample.\allowbreak{}summary()} & \raggedright\arraybackslash\small Return a small text table describing the sample. \\
\small method & \raggedright\arraybackslash\small \texttt{DiscSample.\allowbreak{}with\_\allowbreak{}plane\_\allowbreak{}angle(plane\_\allowbreak{}angle\_\allowbreak{}deg)} & \raggedright\arraybackslash\small Return a copy of this disc with another orientation. \\
\small method & \raggedright\arraybackslash\small \texttt{DiscSample.\allowbreak{}face\_\allowbreak{}area\_\allowbreak{}in\_\allowbreak{}beam(beam\_\allowbreak{}width, beam\_\allowbreak{}height)} & \raggedright\arraybackslash\small Area (mm\textasciicircum{}2) of the part of the beam that hits the FRONT FACE of the disc, seen along the beam (uniform rectangular beam centred on the disc). \\
\small method & \raggedright\arraybackslash\small \texttt{DiscSample.\allowbreak{}contains(points, tol=1e-09)} & \raggedright\arraybackslash\small True for every point (x', y', z') inside the volume of the disc (empty parts of partially filled layers included): distance from the axis <= R and depth between the front and back faces. \\
\small method & \raggedright\arraybackslash\small \texttt{DiscSample.\allowbreak{}\_\allowbreak{}clip\_\allowbreak{}inside(points)} & \raggedright\arraybackslash\small Push points that rounding errors put just outside the disc back onto its surface: depth limited to the faces, and points slightly beyond the rim moved radially onto it. \emph{(internal)} \\
\small method & \raggedright\arraybackslash\small \texttt{DiscSample.\allowbreak{}\_\allowbreak{}section\_\allowbreak{}interval(p, d)} & \raggedright\arraybackslash\small Interval of t during which rays p + t*d are inside the disc (all layers together): inside the cylinder x'\textasciicircum{}2 + y'\textasciicircum{}2 <= R\textasciicircum{}2 AND between the front and back faces. \emph{(internal)} \\
\small function & \raggedright\arraybackslash\small \texttt{create\_\allowbreak{}disc\_\allowbreak{}sample(diameter, layers, z\_\allowbreak{}front=0.\allowbreak{}0, plane\_\allowbreak{}angle\_\allowbreak{}deg=90.\allowbreak{}0, beam\_\allowbreak{}height=None)} & \raggedright\arraybackslash\small Build a layered disc (circular plate) centred on the beam. \\
\bottomrule
\end{longtable}

\section{\texorpdfstring{\texttt{cylinder.py}}{cylinder.py}: Cylindrical geometry}
\label{app:api:cylinder}
CYLINDRICAL geometry of MCAMS: concentric cylinders (an.

\begin{longtable}{@{}p{1.75cm}p{6.35cm}p{6.45cm}@{}}
\toprule \textbf{Kind} & \textbf{Name and arguments} & \textbf{What it does} \\ \midrule \endhead
\small class & \raggedright\arraybackslash\small \texttt{CylindricalSample(LayeredSample)} & \raggedright\arraybackslash\small Concentric vertical cylinders. \\
\small init & \raggedright\arraybackslash\small \texttt{CylindricalSample.\allowbreak{}\_\allowbreak{}\_\allowbreak{}init\_\allowbreak{}\_\allowbreak{}(height, layers, y\_\allowbreak{}bottom=None, beam\_\allowbreak{}height=None)} & \raggedright\arraybackslash\small Constructor: creates the object (see the description of the class). \\
\small method & \raggedright\arraybackslash\small \texttt{CylindricalSample.\allowbreak{}summary()} & \raggedright\arraybackslash\small Return a small text table describing the sample. \\
\small method & \raggedright\arraybackslash\small \texttt{CylindricalSample.\allowbreak{}to\_\allowbreak{}sample\_\allowbreak{}frame(vectors)} & \raggedright\arraybackslash\small Lab -> sample frame. \\
\small method & \raggedright\arraybackslash\small \texttt{CylindricalSample.\allowbreak{}to\_\allowbreak{}lab\_\allowbreak{}frame(vectors)} & \raggedright\arraybackslash\small Sample -> lab frame (identical frames, see above). \\
\small method & \raggedright\arraybackslash\small \texttt{CylindricalSample.\allowbreak{}contains(points, tol=1e-09)} & \raggedright\arraybackslash\small True for every point (x, y, z) inside the cylinder (the empty part of partially filled layers included). \texttt{tol} forgives floating-point rounding errors for points on a surface. \\
\small method & \raggedright\arraybackslash\small \texttt{CylindricalSample.\allowbreak{}layer\_\allowbreak{}of\_\allowbreak{}points(points)} & \raggedright\arraybackslash\small Index of the layer whose MATERIAL contains each point: -1 if the point is outside the sample or in the empty space below or above the material of a partially filled layer. \\
\small static & \raggedright\arraybackslash\small \texttt{CylindricalSample.\allowbreak{}\_\allowbreak{}radius\_\allowbreak{}of(x, y, z)} & \raggedright\arraybackslash\small Distance of the points (x, y, z) to the vertical axis: this number tells in which layer a point is. \emph{(internal)} \\
\small method & \raggedright\arraybackslash\small \texttt{CylindricalSample.\allowbreak{}\_\allowbreak{}shell\_\allowbreak{}interval(p, d, R)} & \raggedright\arraybackslash\small Distances t at which the rays p + t*d enter and leave the (infinitely tall) cylinder of radius R around the vertical axis. \emph{(internal)} \\
\small method & \raggedright\arraybackslash\small \texttt{CylindricalSample.\allowbreak{}\_\allowbreak{}outline\_\allowbreak{}interval(p, d)} & \raggedright\arraybackslash\small Distances t during which the rays are inside the whole sample: inside the outer cylinder AND between the bottom and top faces (heights y\_bottom and y\_top). \emph{(internal)} \\
\small method & \raggedright\arraybackslash\small \texttt{CylindricalSample.\allowbreak{}\_\allowbreak{}trace(points, directions)} & \raggedright\arraybackslash\small Cut the path of each particle inside the sample into segments. \emph{(internal)} \\
\small method & \raggedright\arraybackslash\small \texttt{CylindricalSample.\allowbreak{}outgoing\_\allowbreak{}path\_\allowbreak{}lengths(points, directions)} & \raggedright\arraybackslash\small Distance travelled in the MATERIAL of each layer by particles that start at \texttt{points} (inside the cylinder) and fly along \texttt{directions} until they leave the sample (through the side or through the top/bottom faces). \\
\small method & \raggedright\arraybackslash\small \texttt{CylindricalSample.\allowbreak{}\_\allowbreak{}first\_\allowbreak{}interaction(points, directions, rng)} & \raggedright\arraybackslash\small Find the first interaction of particles flying from \texttt{points} along \texttt{directions} (the points may be inside the sample or upstream of it). \emph{(internal)} \\
\small method & \raggedright\arraybackslash\small \texttt{CylindricalSample.\allowbreak{}next\_\allowbreak{}interaction(points, directions, rng=None)} & \raggedright\arraybackslash\small Next interaction of particles flying inside the sample (used for multiple scattering). \\
\small method & \raggedright\arraybackslash\small \texttt{CylindricalSample.\allowbreak{}sample\_\allowbreak{}interaction\_\allowbreak{}points(x, y, rng=None)} & \raggedright\arraybackslash\small First interaction of beam particles. \\
\small method & \raggedright\arraybackslash\small \texttt{CylindricalSample.\allowbreak{}\_\allowbreak{}check\_\allowbreak{}points(points)} & \raggedright\arraybackslash\small Return points as an (N, 3) float array; raise if outside sample. \emph{(internal)} \\
\small function & \raggedright\arraybackslash\small \texttt{create\_\allowbreak{}cylindrical\_\allowbreak{}sample(height, layers, y\_\allowbreak{}bottom=None, beam\_\allowbreak{}height=None)} & \raggedright\arraybackslash\small Build a sample of concentric vertical cylinders. \\
\bottomrule
\end{longtable}

\section{\texorpdfstring{\texttt{sphere.py}}{sphere.py}: Spherical geometry (version 2.4)}
\label{app:api:sphere}
SPHERICAL geometry of MCAMS (version 2.4): concentric spheres,.

\begin{longtable}{@{}p{1.75cm}p{6.35cm}p{6.45cm}@{}}
\toprule \textbf{Kind} & \textbf{Name and arguments} & \textbf{What it does} \\ \midrule \endhead
\small class & \raggedright\arraybackslash\small \texttt{SphericalSample(CylindricalSample)} & \raggedright\arraybackslash\small Concentric spheres centred on the beam. \\
\small init & \raggedright\arraybackslash\small \texttt{SphericalSample.\allowbreak{}\_\allowbreak{}\_\allowbreak{}init\_\allowbreak{}\_\allowbreak{}(layers, beam\_\allowbreak{}height=None)} & \raggedright\arraybackslash\small Constructor: creates the object (see the description of the class). \\
\small method & \raggedright\arraybackslash\small \texttt{SphericalSample.\allowbreak{}summary()} & \raggedright\arraybackslash\small Return a small text table describing the sample. \\
\small method & \raggedright\arraybackslash\small \texttt{SphericalSample.\allowbreak{}section\_\allowbreak{}area\_\allowbreak{}in\_\allowbreak{}beam(beam\_\allowbreak{}width, beam\_\allowbreak{}height)} & \raggedright\arraybackslash\small Area (mm\textasciicircum{}2) of the part of a uniform rectangular beam (centred on the sphere) that hits the sphere: the circle of radius R\_out seen along the beam, inside the beam rectangle. \\
\small method & \raggedright\arraybackslash\small \texttt{SphericalSample.\allowbreak{}contains(points, tol=1e-09)} & \raggedright\arraybackslash\small True for every point (x, y, z) inside the outer sphere (the empty part of partially filled layers included). \\
\small static & \raggedright\arraybackslash\small \texttt{SphericalSample.\allowbreak{}\_\allowbreak{}radius\_\allowbreak{}of(x, y, z)} & \raggedright\arraybackslash\small Distance of the points to the CENTRE: it tells the layer. \emph{(internal)} \\
\small method & \raggedright\arraybackslash\small \texttt{SphericalSample.\allowbreak{}\_\allowbreak{}shell\_\allowbreak{}interval(p, d, R)} & \raggedright\arraybackslash\small Distances t at which the rays p + t*d enter and leave the sphere of radius R centred on the origin. \emph{(internal)} \\
\small method & \raggedright\arraybackslash\small \texttt{SphericalSample.\allowbreak{}\_\allowbreak{}outline\_\allowbreak{}interval(p, d)} & \raggedright\arraybackslash\small Distances t during which the rays are inside the whole sample: the outer sphere (a sphere has no flat ends). \emph{(internal)} \\
\small function & \raggedright\arraybackslash\small \texttt{create\_\allowbreak{}spherical\_\allowbreak{}sample(layers, beam\_\allowbreak{}height=None)} & \raggedright\arraybackslash\small Build a sample of concentric spheres centred on the beam. \\
\bottomrule
\end{longtable}

\section{\texorpdfstring{\texttt{sample\_factory.py}}{sample_factory.py}: Choice of the geometry}
\label{app:api:sample_factory}
MCAMS: build a sample of the chosen geometry ("slab",.

\begin{longtable}{@{}p{1.75cm}p{6.35cm}p{6.45cm}@{}}
\toprule \textbf{Kind} & \textbf{Name and arguments} & \textbf{What it does} \\ \midrule \endhead
\small function & \raggedright\arraybackslash\small \texttt{make\_\allowbreak{}sample(geometry, layers, height=None, width=None, plane\_\allowbreak{}angle\_\allowbreak{}deg=90.\allowbreak{}0, y\_\allowbreak{}bottom=None, beam\_\allowbreak{}height=None, diameter=None)} & \raggedright\arraybackslash\small Create a sample of the requested geometry. \\
\bottomrule
\end{longtable}

\section{\texorpdfstring{\texttt{transmission\_vs\_angle.py}}{transmission_vs_angle.py}: Transmission versus scattering angle}
\label{app:api:transmission_vs_angle}
Transmission of the scattered particles as a.

\begin{longtable}{@{}p{1.75cm}p{6.35cm}p{6.45cm}@{}}
\toprule \textbf{Kind} & \textbf{Name and arguments} & \textbf{What it does} \\ \midrule \endhead
\small function & \raggedright\arraybackslash\small \texttt{transmission\_\allowbreak{}vs\_\allowbreak{}angle(sample, beam\_\allowbreak{}width, beam\_\allowbreak{}height, H=None, V=None, n\_\allowbreak{}particles=20000, angles\_\allowbreak{}deg=None, phi\_\allowbreak{}deg=0.\allowbreak{}0, rng=None)} & \raggedright\arraybackslash\small Compute the transmission of scattered particles for a list of angles. \\
\small function & \raggedright\arraybackslash\small \texttt{plot\_\allowbreak{}transmission(angles\_\allowbreak{}deg, transmission, title=None, save\_\allowbreak{}as=None, show=True)} & \raggedright\arraybackslash\small Plot the transmission as a function of the scattering angle. \\
\bottomrule
\end{longtable}

\section{\texorpdfstring{\texttt{multiple\_scattering.py}}{multiple_scattering.py}: Single and multiple scattering}
\label{app:api:multiple_scattering}
Count the particles that leave the sample after.

\begin{longtable}{@{}p{1.75cm}p{6.35cm}p{6.45cm}@{}}
\toprule \textbf{Kind} & \textbf{Name and arguments} & \textbf{What it does} \\ \midrule \endhead
\small function & \raggedright\arraybackslash\small \texttt{isotropic\_\allowbreak{}cos(n, rng)} & \raggedright\arraybackslash\small Isotropic scattering: cos(chi) uniform between -1 and 1. (Uniform in cos(chi), not in chi, gives equal probability per unit of solid angle.). \\
\small function & \raggedright\arraybackslash\small \texttt{henyey\_\allowbreak{}greenstein(g)} & \raggedright\arraybackslash\small Return a deflection sampler with the Henyey-Greenstein distribution. \\
\small function & \raggedright\arraybackslash\small \texttt{rotate\_\allowbreak{}directions(d, cos\_\allowbreak{}chi, rng)} & \raggedright\arraybackslash\small New directions that make an angle chi (cos\_chi given) with the unit vectors \texttt{d} (shape (N, 3)), with a random azimuth around \texttt{d}. \\
\small dataclass & \raggedright\arraybackslash\small \texttt{ScatteringResult(angle\_\allowbreak{}edges\_\allowbreak{}deg: np.\allowbreak{}ndarray, angles\_\allowbreak{}deg: np.\allowbreak{}ndarray, solid\_\allowbreak{}angle: np.\allowbreak{}ndarray, counts\_\allowbreak{}by\_\allowbreak{}order: np.\allowbreak{}ndarray, n\_\allowbreak{}incident: int, n\_\allowbreak{}no\_\allowbreak{}interaction: int, n\_\allowbreak{}absorbed: int, n\_\allowbreak{}exited: int, n\_\allowbreak{}unfinished: int)} & \raggedright\arraybackslash\small Result of \texttt{simulate\_\allowbreak{}scattering\_\allowbreak{}orders}. \\
\small property & \raggedright\arraybackslash\small \texttt{ScatteringResult.\allowbreak{}single} & \raggedright\arraybackslash\small Counts per bin of particles scattered exactly once. \\
\small property & \raggedright\arraybackslash\small \texttt{ScatteringResult.\allowbreak{}multiple} & \raggedright\arraybackslash\small Counts per bin of particles scattered two or more times. \\
\small property & \raggedright\arraybackslash\small \texttt{ScatteringResult.\allowbreak{}total} & \raggedright\arraybackslash\small Counts per bin of all scattered particles (single + multiple). \\
\small method & \raggedright\arraybackslash\small \texttt{ScatteringResult.\allowbreak{}per\_\allowbreak{}steradian(counts)} & \raggedright\arraybackslash\small Convert counts per bin into particles per incident particle per steradian, which does not depend on the bin width. \\
\small function & \raggedright\arraybackslash\small \texttt{simulate\_\allowbreak{}scattering\_\allowbreak{}orders(sample, beam\_\allowbreak{}width, beam\_\allowbreak{}height, H=None, V=None, n\_\allowbreak{}particles=200000, angle\_\allowbreak{}edges\_\allowbreak{}deg=None, phi\_\allowbreak{}deg=0.\allowbreak{}0, phi\_\allowbreak{}halfwidth\_\allowbreak{}deg=None, deflection\_\allowbreak{}sampler=None, max\_\allowbreak{}orders=200, rng=None)} & \raggedright\arraybackslash\small Simulate particles through the sample and count them by scattering angle and by number of scatterings (single / multiple). \\
\small function & \raggedright\arraybackslash\small \texttt{plot\_\allowbreak{}scattering\_\allowbreak{}orders(result, per\_\allowbreak{}steradian=True, logy=False, title=None, save\_\allowbreak{}as=None, show=True)} & \raggedright\arraybackslash\small Plot total, single and multiple scattering as a function of 2theta. \\
\bottomrule
\end{longtable}

\section{\texorpdfstring{\texttt{transmission\_by\_layer.py}}{transmission_by_layer.py}: Transmission by layer of last scattering}
\label{app:api:transmission_by_layer}
Transmission vs scattering angle, separated by the.

\begin{longtable}{@{}p{1.75cm}p{6.35cm}p{6.45cm}@{}}
\toprule \textbf{Kind} & \textbf{Name and arguments} & \textbf{What it does} \\ \midrule \endhead
\small dataclass & \raggedright\arraybackslash\small \texttt{LastLayerResult(angles\_\allowbreak{}deg: np.\allowbreak{}ndarray, layer\_\allowbreak{}names: tuple, sum\_\allowbreak{}attenuation: np.\allowbreak{}ndarray, n\_\allowbreak{}events: np.\allowbreak{}ndarray, n\_\allowbreak{}incident: int)} & \raggedright\arraybackslash\small Result of \texttt{transmission\_\allowbreak{}by\_\allowbreak{}last\_\allowbreak{}layer}. \\
\small static & \raggedright\arraybackslash\small \texttt{LastLayerResult.\allowbreak{}\_\allowbreak{}classes(kind)} & \raggedright\arraybackslash\small Which order classes belong to \texttt{kind} ("single", "multiple", "total"). \emph{(internal)} \\
\small method & \raggedright\arraybackslash\small \texttt{LastLayerResult.\allowbreak{}events(kind='total')} & \raggedright\arraybackslash\small Number of scattering events in each layer, shape (L,). \\
\small method & \raggedright\arraybackslash\small \texttt{LastLayerResult.\allowbreak{}transmission\_\allowbreak{}of(kind='total')} & \raggedright\arraybackslash\small Transmission vs angle for every layer of last scattering, shape (L, A). \\
\small method & \raggedright\arraybackslash\small \texttt{LastLayerResult.\allowbreak{}mean\_\allowbreak{}transmission(kind='total')} & \raggedright\arraybackslash\small Transmission vs angle averaged over ALL the layers, weighted by the number of events of each layer, shape (A,). \\
\small property & \raggedright\arraybackslash\small \texttt{LastLayerResult.\allowbreak{}transmission} & \raggedright\arraybackslash\small (L, A): last scattering in layer k, all scattering orders. \\
\small property & \raggedright\arraybackslash\small \texttt{LastLayerResult.\allowbreak{}transmission\_\allowbreak{}single} & \raggedright\arraybackslash\small (L, A): last scattering in layer k, particles scattered once. \\
\small property & \raggedright\arraybackslash\small \texttt{LastLayerResult.\allowbreak{}transmission\_\allowbreak{}multiple} & \raggedright\arraybackslash\small (L, A): last scattering in layer k, particles scattered 2+ times. \\
\small function & \raggedright\arraybackslash\small \texttt{transmission\_\allowbreak{}by\_\allowbreak{}last\_\allowbreak{}layer(sample, beam\_\allowbreak{}width, beam\_\allowbreak{}height, H=None, V=None, n\_\allowbreak{}particles=50000, angles\_\allowbreak{}deg=None, phi\_\allowbreak{}deg=0.\allowbreak{}0, deflection\_\allowbreak{}sampler=None, max\_\allowbreak{}orders=200, rng=None)} & \raggedright\arraybackslash\small Transmission vs scattering angle, with one vector per layer in which the last scattering took place (see the top of this file). \\
\small function & \raggedright\arraybackslash\small \texttt{plot\_\allowbreak{}transmission\_\allowbreak{}by\_\allowbreak{}layer(result, kind='all', title=None, save\_\allowbreak{}as=None, show=True)} & \raggedright\arraybackslash\small Plot the transmission vs angle, one curve per layer of last scattering. \\
\bottomrule
\end{longtable}

\section{\texorpdfstring{\texttt{direct\_attenuation.py}}{direct_attenuation.py}: Direct method for single-scattering attenuation}
\label{app:api:direct_attenuation}
SINGLE-SCATTERING attenuation factors computed by.

\begin{longtable}{@{}p{1.75cm}p{6.35cm}p{6.45cm}@{}}
\toprule \textbf{Kind} & \textbf{Name and arguments} & \textbf{What it does} \\ \midrule \endhead
\small dataclass & \raggedright\arraybackslash\small \texttt{AttenuationResult(angles\_\allowbreak{}deg: np.\allowbreak{}ndarray, layer\_\allowbreak{}names: tuple, factors: np.\allowbreak{}ndarray, errors: np.\allowbreak{}ndarray, mean\_\allowbreak{}path: np.\allowbreak{}ndarray, mean\_\allowbreak{}path\_\allowbreak{}error: np.\allowbreak{}ndarray, mus: np.\allowbreak{}ndarray, scatter\_\allowbreak{}fractions: np.\allowbreak{}ndarray, method: str, n\_\allowbreak{}samples: np.\allowbreak{}ndarray = field(default\_\allowbreak{}factory=lambda: np.\allowbreak{}zeros(0)), seconds: float = 0.\allowbreak{}0)} & \raggedright\arraybackslash\small Single-scattering attenuation factors, one row per layer. \\
\small method & \raggedright\arraybackslash\small \texttt{AttenuationResult.\allowbreak{}layer(name)} & \raggedright\arraybackslash\small Index of the layer called \texttt{name}. \\
\small method & \raggedright\arraybackslash\small \texttt{AttenuationResult.\allowbreak{}single\_\allowbreak{}per\_\allowbreak{}steradian(isotropic=True)} & \raggedright\arraybackslash\small Singly scattered particles per incident particle and per steradian that leave the sample towards each angle (sum over the layers):. \\
\small function & \raggedright\arraybackslash\small \texttt{\_\allowbreak{}lab\_\allowbreak{}z\_\allowbreak{}range(sample)} & \raggedright\arraybackslash\small Interval of lab z (beam direction) that contains the whole sample. \emph{(internal)} \\
\small function & \raggedright\arraybackslash\small \texttt{\_\allowbreak{}draw\_\allowbreak{}points(sample, beam\_\allowbreak{}width, beam\_\allowbreak{}height, H, V, n\_\allowbreak{}target, max\_\allowbreak{}draws, chunk, rng)} & \raggedright\arraybackslash\small Draw points uniformly in the volume lit by the beam and keep, for every layer, up to \texttt{n\_\allowbreak{}target} of those that fall in its MATERIAL. \emph{(internal)} \\
\small function & \raggedright\arraybackslash\small \texttt{\_\allowbreak{}chord\_\allowbreak{}ratio(sample, pts)} & \raggedright\arraybackslash\small For points in the material, the FULL path of their beam line in each layer (backwards along -z plus forwards along +z, lab frame), divided by the total path in all layers: array (N, L), each row sums to 1. \emph{(internal)} \\
\small function & \raggedright\arraybackslash\small \texttt{\_\allowbreak{}mean\_\allowbreak{}paths\_\allowbreak{}from\_\allowbreak{}draws(sample, counts, n\_\allowbreak{}draws, zlen, pts\_\allowbreak{}uniform)} & \raggedright\arraybackslash\small Mean path <L\_k> of the beam in each layer, with its error. \emph{(internal)} \\
\small function & \raggedright\arraybackslash\small \texttt{mean\_\allowbreak{}material\_\allowbreak{}paths(sample, beam\_\allowbreak{}width, beam\_\allowbreak{}height, H=None, V=None, n\_\allowbreak{}draws=1000000, n\_\allowbreak{}lines=200000, rng=None)} & \raggedright\arraybackslash\small Mean path (mm) of the incident beam in the material of each layer, averaged over the whole beam (H, V profiles, including the parts of the beam that miss the layer), with its statistical error. \\
\small function & \raggedright\arraybackslash\small \texttt{attenuation\_\allowbreak{}factors\_\allowbreak{}direct(sample, beam\_\allowbreak{}width, beam\_\allowbreak{}height, H=None, V=None, angles\_\allowbreak{}deg=None, phi\_\allowbreak{}deg=0.\allowbreak{}0, n\_\allowbreak{}points=200000, max\_\allowbreak{}draws=None, chunk=500000, rng=None)} & \raggedright\arraybackslash\small Single-scattering attenuation factors A\_k(2theta) by direct sampling. \\
\small function & \raggedright\arraybackslash\small \texttt{attenuation\_\allowbreak{}factors\_\allowbreak{}mc(sample, beam\_\allowbreak{}width, beam\_\allowbreak{}height, H=None, V=None, angles\_\allowbreak{}deg=None, phi\_\allowbreak{}deg=0.\allowbreak{}0, n\_\allowbreak{}particles=200000, n\_\allowbreak{}batches=5, mean\_\allowbreak{}path=None, rng=None)} & \raggedright\arraybackslash\small The same A\_k(2theta), computed with the analog Monte Carlo of MCAMS (transmission\_by\_layer.transmission\_by\_last\_layer, first scatterings only), for comparison with the direct method. \\
\small function & \raggedright\arraybackslash\small \texttt{compare\_\allowbreak{}results(reference, other, every=30)} & \raggedright\arraybackslash\small Text table comparing two AttenuationResult (e.g. direct and MC): for each layer, the values every \texttt{every} degrees, the relative difference and the difference in units of the combined error (z). |z| < 2 for most angles means agreement within the statistics. \\
\small function & \raggedright\arraybackslash\small \texttt{plot\_\allowbreak{}attenuation\_\allowbreak{}factors(results, title=None, save\_\allowbreak{}as=None, show=True)} & \raggedright\arraybackslash\small Plot A\_k(2theta) of one or several AttenuationResult. \\
\bottomrule
\end{longtable}

\section{\texorpdfstring{\texttt{formula.py}}{formula.py}: Chemical formulas}
\label{app:api:formula}
Read a chemical formula such as "H2O", "Fe2O3Na2Li" or "[2H]2O".

\begin{longtable}{@{}p{1.75cm}p{6.35cm}p{6.45cm}@{}}
\toprule \textbf{Kind} & \textbf{Name and arguments} & \textbf{What it does} \\ \midrule \endhead
\small class & \raggedright\arraybackslash\small \texttt{FormulaError(ValueError)} & \raggedright\arraybackslash\small The chemical formula (or a species name) is not valid. \\
\small dataclass & \raggedright\arraybackslash\small \texttt{Species(symbol: str, mass\_\allowbreak{}number: Optional[int] = None)} & \raggedright\arraybackslash\small A kind of atom. \\
\small property & \raggedright\arraybackslash\small \texttt{Species.\allowbreak{}key} & \raggedright\arraybackslash\small Short text without brackets: "Fe" for natural iron, "56Fe" for the isotope. \\
\small function & \raggedright\arraybackslash\small \texttt{\_\allowbreak{}make\_\allowbreak{}species(symbol, mass\_\allowbreak{}number=None)} & \raggedright\arraybackslash\small Create a Species after checking the symbol (and the D/T shorthands). \emph{(internal)} \\
\small function & \raggedright\arraybackslash\small \texttt{parse\_\allowbreak{}species(text)} & \raggedright\arraybackslash\small Convert the text of one species ("Fe", "56Fe", "[56Fe]", "D") into a \texttt{Species}. \\
\small function & \raggedright\arraybackslash\small \texttt{\_\allowbreak{}add(counts, species, amount)} & \raggedright\arraybackslash\small Add \texttt{amount} atoms of \texttt{species} to the dictionary \texttt{counts}. \emph{(internal)} \\
\small function & \raggedright\arraybackslash\small \texttt{\_\allowbreak{}parse\_\allowbreak{}part(text, original)} & \raggedright\arraybackslash\small Parse a formula without hydrate separators. \emph{(internal)} \\
\small function & \raggedright\arraybackslash\small \texttt{parse\_\allowbreak{}formula(formula)} & \raggedright\arraybackslash\small Convert a chemical formula into \{Species: number of atoms\}. \\
\small function & \raggedright\arraybackslash\small \texttt{format\_\allowbreak{}composition(composition)} & \raggedright\arraybackslash\small Readable text of a composition, e.g. "[2H]:2 O:1". \\
\bottomrule
\end{longtable}

\section{\texorpdfstring{\texttt{nuclear\_data.py}}{nuclear_data.py}: Atomic and nuclear data}
\label{app:api:nuclear_data}
Atomic masses and neutron cross sections of the elements and.

\begin{longtable}{@{}p{1.75cm}p{6.35cm}p{6.45cm}@{}}
\toprule \textbf{Kind} & \textbf{Name and arguments} & \textbf{What it does} \\ \midrule \endhead
\small class & \raggedright\arraybackslash\small \texttt{TableError(LookupError)} & \raggedright\arraybackslash\small A value is not available in the tables. \\
\small class & \raggedright\arraybackslash\small \texttt{AtomicData} & \raggedright\arraybackslash\small Access to atomic masses and neutron cross sections. \\
\small init & \raggedright\arraybackslash\small \texttt{AtomicData.\allowbreak{}\_\allowbreak{}\_\allowbreak{}init\_\allowbreak{}\_\allowbreak{}(use\_\allowbreak{}periodictable=True)} & \raggedright\arraybackslash\small Constructor: creates the object (see the description of the class). \\
\small property & \raggedright\arraybackslash\small \texttt{AtomicData.\allowbreak{}has\_\allowbreak{}neutron\_\allowbreak{}data} & \raggedright\arraybackslash\small True if tabulated neutron cross sections can be used. \\
\small method & \raggedright\arraybackslash\small \texttt{AtomicData.\allowbreak{}\_\allowbreak{}pt\_\allowbreak{}item(species)} & \raggedright\arraybackslash\small The element (natural) or isotope object of \texttt{periodictable}. \emph{(internal)} \\
\small method & \raggedright\arraybackslash\small \texttt{AtomicData.\allowbreak{}mass(species)} & \raggedright\arraybackslash\small Mass in g/mol of one atom (mole of atoms) of \texttt{species}. \\
\small method & \raggedright\arraybackslash\small \texttt{AtomicData.\allowbreak{}neutron\_\allowbreak{}xs(species)} & \raggedright\arraybackslash\small Tabulated neutron cross sections of ONE atom of \texttt{species}, at the thermal velocity 2200 m/s, in barn. \\
\bottomrule
\end{longtable}

\section{\texorpdfstring{\texttt{layer\_builder.py}}{layer_builder.py}: Layers from formula, density and cross sections}
\label{app:api:layer_builder}
Prepare the input of a Layer from the chemical formula,.

\begin{longtable}{@{}p{1.75cm}p{6.35cm}p{6.45cm}@{}}
\toprule \textbf{Kind} & \textbf{Name and arguments} & \textbf{What it does} \\ \midrule \endhead
\small dataclass & \raggedright\arraybackslash\small \texttt{LayerProperties(label: str, thickness\_\allowbreak{}mm: float, formula: str, density\_\allowbreak{}g\_\allowbreak{}cm3: float, composition: dict, molar\_\allowbreak{}mass: float, formula\_\allowbreak{}units\_\allowbreak{}per\_\allowbreak{}cm3: float, sigma\_\allowbreak{}scat\_\allowbreak{}barn: float, sigma\_\allowbreak{}abs\_\allowbreak{}barn: float, mu\_\allowbreak{}scat: float, mu\_\allowbreak{}abs: float, mu: float, scatter\_\allowbreak{}fraction: float, sources: str)} & \raggedright\arraybackslash\small Everything computed for one layer. \\
\small method & \raggedright\arraybackslash\small \texttt{LayerProperties.\allowbreak{}as\_\allowbreak{}tuple(fill\_\allowbreak{}height=None)} & \raggedright\arraybackslash\small (label, thickness in mm, mu in 1/mm, scatter fraction): the input of create\_slab\_sample / create\_cylindrical\_sample / make\_sample. \\
\small method & \raggedright\arraybackslash\small \texttt{LayerProperties.\allowbreak{}as\_\allowbreak{}layer(fill\_\allowbreak{}height=None, fullness=None, bottom=None, height=None)} & \raggedright\arraybackslash\small The same data as a \texttt{Layer} object, optionally partially filled: give \texttt{bottom} and/or \texttt{height} (mm, from the centre of the beam; version 2.2, recommended), OR \texttt{fill\_\allowbreak{}height} (mm from the bottom of the sample) OR \texttt{fullness} (fraction of the beam height covered ... \\
\small method & \raggedright\arraybackslash\small \texttt{LayerProperties.\allowbreak{}summary()} & \raggedright\arraybackslash\small Readable description of the calculation. \\
\small function & \raggedright\arraybackslash\small \texttt{\_\allowbreak{}species\_\allowbreak{}dict(user\_\allowbreak{}dict)} & \raggedright\arraybackslash\small Convert \{"Fe": value, "56Fe": value, ...\} into \{Species: value\}. \emph{(internal)} \\
\small function & \raggedright\arraybackslash\small \texttt{compute\_\allowbreak{}layer\_\allowbreak{}properties(label, thickness, formula, density, sigma\_\allowbreak{}scat=None, sigma\_\allowbreak{}abs=None, cross\_\allowbreak{}sections=None, masses=None, wavelength\_\allowbreak{}A=None, data=None)} & \raggedright\arraybackslash\small Compute the attenuation coefficient and the scatter fraction of a layer. \\
\small function & \raggedright\arraybackslash\small \texttt{make\_\allowbreak{}layer\_\allowbreak{}tuple(label, thickness, formula, density, sigma\_\allowbreak{}scat=None, sigma\_\allowbreak{}abs=None, cross\_\allowbreak{}sections=None, masses=None, wavelength\_\allowbreak{}A=None, data=None, fill\_\allowbreak{}height=None)} & \raggedright\arraybackslash\small Same as \texttt{compute\_\allowbreak{}layer\_\allowbreak{}properties} (same parameters), but returns only the tuple (label, thickness\_mm, mu\_per\_mm, scatter\_fraction) that defines a Layer, ready to be put in the list of layers of create\_slab\_sample, create\_cylindrical\_sample or make\_sample. \\
\small function & \raggedright\arraybackslash\small \texttt{build\_\allowbreak{}layer(label, thickness, formula, density, sigma\_\allowbreak{}scat=None, sigma\_\allowbreak{}abs=None, cross\_\allowbreak{}sections=None, masses=None, wavelength\_\allowbreak{}A=None, data=None, fill\_\allowbreak{}height=None, fullness=None, bottom=None, height=None)} & \raggedright\arraybackslash\small Same as \texttt{make\_\allowbreak{}layer\_\allowbreak{}tuple}, but returns a \texttt{Layer} object, which can also have a limited vertical extent: \texttt{bottom} and \texttt{height} (mm, from the centre of the beam; version 2.2), or \texttt{fill\_\allowbreak{}height}, or \texttt{fullness} (fraction of the beam height covered by the ... \\
\bottomrule
\end{longtable}

\section{\texorpdfstring{\texttt{run\_simulation.py}}{run_simulation.py}: Main script}
\label{app:api:run_simulation}
Main script of MCAMS (Monte Carlo Attenuation and.

\begin{longtable}{@{}p{1.75cm}p{6.35cm}p{6.45cm}@{}}
\toprule \textbf{Kind} & \textbf{Name and arguments} & \textbf{What it does} \\ \midrule \endhead
\small function & \raggedright\arraybackslash\small \texttt{\_\allowbreak{}with\_\allowbreak{}extent(layers, names, fullness=None, bottom=None, height=None)} & \raggedright\arraybackslash\small Copy of the list of layer tuples in which the layers called \texttt{names} get a limited vertical extent: \texttt{bottom}/\texttt{height} (version 2.2) or \texttt{fullness} (version 2.0). \emph{(internal)} \\
\small function & \raggedright\arraybackslash\small \texttt{build\_\allowbreak{}sample(geometry, plane\_\allowbreak{}angle\_\allowbreak{}deg=90.\allowbreak{}0, fullness=None, bottom=None, height=None)} & \raggedright\arraybackslash\small Create the sample of the chosen geometry from the settings above. \\
\small function & \raggedright\arraybackslash\small \texttt{run\_\allowbreak{}simulation(geometry=DEFAULT\_\allowbreak{}GEOMETRY, plane\_\allowbreak{}angle\_\allowbreak{}deg=90.\allowbreak{}0, n\_\allowbreak{}particles=500000, show=True, output\_\allowbreak{}prefix='', fullness=None, bottom=None, height=None)} & \raggedright\arraybackslash\small Run the complete chain for the chosen geometry. \\
\bottomrule
\end{longtable}
